%% decision-hierarchy-apparatus-supplement.tex
%% Separately authored apparatus following the exact hierarchy chapter.

\section*{Arithmetical-hierarchy apparatus}
\noindent\textit{Reader note.}
This separately authored table and figure summarize the exact inherited
arithmetical-hierarchy chapter.  Their labels and evidence point only to
genuine source sections or results and do not replace the inherited proofs.

\begin{table}[htbp]
\centering
\scriptsize
\begin{tabular}{@{}p{.12\linewidth}p{.27\linewidth}p{.24\linewidth}p{.27\linewidth}@{}}
\toprule
Class & Normal form and closure & Computability reading & Calibrating set/result \\
\midrule
\(\Sigma^0_0=\Pi^0_0\) & bounded negation-normal-form matrix; closed under
Boolean operations and marked bounded quantifiers & bounded-formula level & no
completeness claim made here \\
\(\boldsymbol\Sigma^0_1\) & existential prefix; closed under
\(\land,\lor\), marked bounded quantifiers, and existential quantification &
c.e. sets & \(K\) is many-one complete \\
\(\boldsymbol\Pi^0_1\) & universal-prefix dual of the preceding row &
complements of c.e. sets & \(\N\setminus K\) is many-one complete \\
\(\boldsymbol\Delta^0_1\) & both level-one forms & decidable sets &
the level-one classification theorem in Section~\ref{sec:decision-hierarchy-02} \\
\(\boldsymbol\Pi^0_2\) & \(\forall\vec u\exists\vec v\,\psi\), bounded
\(\psi\) & universal--existential level & \(\mathsf{TOT}\) is many-one
complete \\
\(\boldsymbol\Sigma^0_{n+1}\) & \(n+1\) alternating blocks, existential
first; the closure clauses are those of
Theorem~\ref{thm:prefix-normal-form} & c.e. relative to
\(\varnothing^{(n)}\) & \(\varnothing^{(n+1)}\) is the standard complete
representative \cite{Shoenfield1967} \\
\(\boldsymbol\Pi^0_{n+1}\) & universal-prefix dual & complements of sets c.e.
relative to \(\varnothing^{(n)}\) & the complement of
\(\varnothing^{(n+1)}\) is complete \cite{Shoenfield1967} \\
\(\boldsymbol\Delta^0_{n+1}\) & both level-\(n+1\) forms & exactly the sets
Turing reducible to \(\varnothing^{(n)}\) &
Post's theorem, Theorem~\ref{thm:post-hierarchy}, for \(n\geq0\) \\
\bottomrule
\end{tabular}
\caption{Arithmetical hierarchy.  Formula classes and
boldface relation classes remain distinct; the prefix/closure equivalence is
Theorem~\ref{thm:prefix-normal-form}, and every completeness claim is
under many-one reductions.}
\label{tab:p04-arithmetical-hierarchy}
\end{table}

\begin{figure}[htbp]
\centering
\begin{tikzpicture}[
  node distance=10mm and 16mm,
  cls/.style={draw,rounded corners,minimum width=31mm,minimum height=8mm,
              align=center,font=\small},
  strict/.style={-{Latex[length=2mm]},very thick}
]
 \node[cls] (d1) {\(\boldsymbol\Delta^0_1\)\\decidable};
 \node[cls,above left=of d1] (s1)
   {\(\boldsymbol\Sigma^0_1\)\\\(K\) complete};
 \node[cls,above right=of d1] (p1)
   {\(\boldsymbol\Pi^0_1\)\\\(\N\setminus K\) complete};
 \node[cls,above=25mm of d1] (d2)
   {\(\boldsymbol\Delta^0_2
      =\boldsymbol\Sigma^0_2\cap\boldsymbol\Pi^0_2\)};
 \node[cls,above left=of d2] (s2)
   {\(\boldsymbol\Sigma^0_2\)\\\(\N\setminus\mathsf{TOT}\) complete};
 \node[cls,above right=of d2] (p2)
   {\(\boldsymbol\Pi^0_2\)\\\(\mathsf{TOT}\) complete};
 \node[cls,above=25mm of d2] (d3)
   {\(\boldsymbol\Delta^0_3
      =\boldsymbol\Sigma^0_3\cap\boldsymbol\Pi^0_3\)};
 \draw[strict] (d1) -- (s1);
 \draw[strict] (d1) -- (p1);
 \draw[strict] (s1) -- (d2);
 \draw[strict] (p1) -- (d2);
 \draw[strict] (d2) -- (s2);
 \draw[strict] (d2) -- (p2);
 \draw[strict] (s2) -- (d3);
 \draw[strict] (p2) -- (d3);
 \node[font=\small,above=3mm of d3] {\(\vdots\)};
\end{tikzpicture}
\caption{The first levels of the arithmetical hierarchy.  Every heavy arrow
is a strict inclusion from
Theorem~\ref{thm:arithmetical-hierarchy-strictness}; the intersections
are definitional.  Completeness is under many-one reductions, with
\(K\), its complement, and \(\mathsf{TOT}\) supplied by the inherited
hierarchy-classification theorem,
Theorem~\ref{thm:arithmetical-hierarchy-classification}.  Higher levels repeat
the same pattern with iterated jumps as calibrated by
Theorem~\ref{thm:post-hierarchy}.}
\label{fig:p04-arithmetical-hierarchy}
\end{figure}
